% GENERATED FILE. Edit canonical lessons, then run npm run generate:books.
% canonical-source-hash: fnv1a64:73156b9a9efeab0d
% canonical-lessons: PT-C331-temperar, PT-C331-tempero, PT-C331-refogar, PT-C331-refogado, PT-C331-lume-brando

\chapter{Cebola e alho}
\label{ch:pt-cebola-e-alho}
\hlchaptermodality{331}

\begin{chapteropening}
\textbf{By the end of this chapter:} \emph{I can describe seasoning food and starting a dish with a gentle fry on a low heat.}

Complete the last lesson of chapter 331: I can describe seasoning food and starting a dish with a gentle fry on a low heat.
\end{chapteropening}

\section[temperar]{temperar — to season}

\label{lesson:PT-C331-temperar}



\begin{quote}
Before the new one: say the Portuguese for to shred, then the Portuguese for the bowl.
\end{quote}

\subsection*{You'll want to know: temperar}
\textbf{temperar} — \textquotedblleft{}to season\textquotedblright{}.

\begin{quote}
A avó tempera a carne com alho e sal.
\end{quote}

\begin{quote}
\emph{Grandmother seasons the meat with garlic and salt.}
\end{quote}

\emph{Temperar a salada} — to dress the salad.

\begin{cousinweb}
\emph{temperar} comes from Latin \emph{temperāre} \textquotedblleft{}to mix in due proportion\textquotedblright{} (EtymDB), related to \emph{tempus} \textquotedblleft{}time, the right moment\textquotedblright{}.
\end{cousinweb}

\begin{grammarlens}[title={What you've built}]
One more word for telling, explaining and describing at B1.
\end{grammarlens}

\subsection*{Your turn}
\begin{itemize}
\raggedright
  \item \emph{Say it:} \emph{temperar}
  \item \emph{Say it:} \emph{temperar}, once more
  \item \emph{Say it:} a sentence of your own with \emph{temperar}
  \item \emph{Recall:} say the Portuguese for to shred, then the Portuguese for the bowl, then say \emph{temperar} again
\end{itemize}

\begin{tcolorbox}[breakable,colback=teal!4,colframe=teal!35!black,title={Before you move on}]
What does temperar mean? (\textquotedblleft{}To season\textquotedblright{}.) Say it once more.
\end{tcolorbox}
\section[o tempero]{o tempero — the seasoning}

\label{lesson:PT-C331-tempero}



\begin{quote}
Before the new one: say the Portuguese for the bowl, then the Portuguese for to season.
\end{quote}

\subsection*{You'll want to know: o tempero}
\textbf{o tempero} — \textquotedblleft{}the seasoning\textquotedblright{}.

\begin{quote}
O segredo deste prato é o tempero.
\end{quote}

\begin{quote}
\emph{The secret of this dish is the seasoning.}
\end{quote}

Built on \emph{temperar}.

\begin{cousinweb}
\emph{tempero} is built on \emph{temperar} \textquotedblleft{}to season\textquotedblright{} (EtymDB), from Latin \emph{temperāre} \textquotedblleft{}to mix in due proportion\textquotedblright{}.
\end{cousinweb}

\begin{grammarlens}[title={What you've built}]
One more word for telling, explaining and describing at B1.
\end{grammarlens}

\subsection*{Your turn}
\begin{itemize}
\raggedright
  \item \emph{Say it:} \emph{o tempero}
  \item \emph{Say it:} \emph{o tempero}, once more
  \item \emph{Say it:} a sentence of your own with \emph{o tempero}
  \item \emph{Recall:} say the Portuguese for the bowl, then the Portuguese for to season, then say \emph{o tempero} again
\end{itemize}

\begin{tcolorbox}[breakable,colback=teal!4,colframe=teal!35!black,title={Before you move on}]
What does o tempero mean? (\textquotedblleft{}The seasoning\textquotedblright{}.) Say it once more.
\end{tcolorbox}
\section[refogar]{refogar — to sauté gently, to soften in oil}

\label{lesson:PT-C331-refogar}



\begin{quote}
Before the new one: say the Portuguese for to season, then the Portuguese for the seasoning.
\end{quote}

\subsection*{You'll want to know: refogar}
\textbf{refogar} — \textquotedblleft{}to sauté gently, to soften in oil\textquotedblright{}.

\begin{quote}
Primeiro, o cozinheiro refoga a cebola em azeite.
\end{quote}

\begin{quote}
\emph{First the cook softens the onion in olive oil.}
\end{quote}

The first step of countless Portuguese dishes.

\begin{cousinweb}
\emph{refogar} is usually traced to \emph{re-} and \emph{fogo} \textquotedblleft{}fire\textquotedblright{}, from Latin \emph{focus} \textquotedblleft{}a hearth\textquotedblright{} (Lewis \& Short); EtymDB has no entry.
\end{cousinweb}

\begin{grammarlens}[title={What you've built}]
One more word for telling, explaining and describing at B1.
\end{grammarlens}

\subsection*{Your turn}
\begin{itemize}
\raggedright
  \item \emph{Say it:} \emph{refogar}
  \item \emph{Say it:} \emph{refogar}, once more
  \item \emph{Say it:} a sentence of your own with \emph{refogar}
  \item \emph{Recall:} say the Portuguese for to season, then the Portuguese for the seasoning, then say \emph{refogar} again
\end{itemize}

\begin{tcolorbox}[breakable,colback=teal!4,colframe=teal!35!black,title={Before you move on}]
What does refogar mean? (\textquotedblleft{}To sauté gently, to soften in oil\textquotedblright{}.) Say it once more.
\end{tcolorbox}
\section[o refogado]{o refogado — the onion and olive oil base}

\label{lesson:PT-C331-refogado}



\begin{quote}
Before the new one: say the Portuguese for the seasoning, then the Portuguese for to sauté gently.
\end{quote}

\subsection*{You'll want to know: o refogado}
\textbf{o refogado} — \textquotedblleft{}the onion and olive oil base\textquotedblright{}.

\begin{quote}
O refogado é a base de muitos pratos portugueses.
\end{quote}

\begin{quote}
\emph{The onion and olive oil base is the start of many Portuguese dishes.}
\end{quote}

Onion, garlic and olive oil, sometimes with tomato.

\begin{cousinweb}
\emph{refogado} is the past participle of \emph{refogar} used as a noun; the verb is usually traced to \emph{fogo} \textquotedblleft{}fire\textquotedblright{}, from Latin \emph{focus} (Lewis \& Short).
\end{cousinweb}

\begin{grammarlens}[title={What you've built}]
One more word for telling, explaining and describing at B1.
\end{grammarlens}

\subsection*{Your turn}
\begin{itemize}
\raggedright
  \item \emph{Say it:} \emph{o refogado}
  \item \emph{Say it:} \emph{o refogado}, once more
  \item \emph{Say it:} a sentence of your own with \emph{o refogado}
  \item \emph{Recall:} say the Portuguese for the seasoning, then the Portuguese for to sauté gently, then say \emph{o refogado} again
\end{itemize}

\begin{tcolorbox}[breakable,colback=teal!4,colframe=teal!35!black,title={Before you move on}]
What does o refogado mean? (\textquotedblleft{}The onion and olive oil base\textquotedblright{}.) Say it once more.
\end{tcolorbox}
\section[o lume brando]{o lume brando — low heat, a low flame}

\label{lesson:PT-C331-lume-brando}



\begin{quote}
Before the new one: say the Portuguese for to sauté gently, then the Portuguese for the onion and olive oil base.
\end{quote}

\subsection*{You'll want to know: o lume brando}
\textbf{o lume brando} — \textquotedblleft{}low heat, a low flame\textquotedblright{}.

\begin{quote}
A sopa ficou ao lume brando durante uma hora.
\end{quote}

\begin{quote}
\emph{The soup stayed on a low heat for an hour.}
\end{quote}

\emph{Em lume brando} — on a low heat. In Brazil: \emph{fogo baixo}.

\begin{cousinweb}
\emph{lume} goes back to Latin \emph{lūmen} \textquotedblleft{}light\textquotedblright{}, and \emph{brando} to Latin \emph{blandus} \textquotedblleft{}gentle, soft\textquotedblright{} (EtymDB).
\end{cousinweb}

\begin{grammarlens}[title={What you've built}]
One more word for telling, explaining and describing at B1.
\end{grammarlens}

\subsection*{Your turn}
\begin{itemize}
\raggedright
  \item \emph{Say it:} \emph{o lume brando}
  \item \emph{Say it:} \emph{o lume brando}, once more
  \item \emph{Say it:} a sentence of your own with \emph{o lume brando}
  \item \emph{Recall:} say the Portuguese for to sauté gently, then the Portuguese for the onion and olive oil base, then say \emph{o lume brando} again
\end{itemize}

\begin{tcolorbox}[breakable,colback=teal!4,colframe=teal!35!black,title={Before you move on}]
What does o lume brando mean? (\textquotedblleft{}Low heat, a low flame\textquotedblright{}.) Say it once more.
\end{tcolorbox}
